\documentclass[11pt]{article}
\usepackage[review]{acl}
\usepackage{times}
\usepackage{latexsym}
\usepackage[T1]{fontenc}
\usepackage[utf8]{inputenc}
\usepackage{microtype}
\usepackage{inconsolata}
\usepackage{lipsum}
\usepackage{url}

\title{Modern LLM Techniques: arXiv-Style References}

\author{First Author \\
  Test University \\
  \texttt{first@test.edu} \\\And
  Second Author \\
  Test University \\
  \texttt{second@test.edu} \\}

\begin{document}
\maketitle

\begin{abstract}
This paper tests extraction of references in informal arXiv preprint style:
no venue name, no page numbers, ``et al.'' author truncation, and DOI/URL fields.
These patterns are common in fast-moving ML subfields where authors cite
preprints before they are formally published.
\end{abstract}

\section{Introduction}
\lipsum[1-2]
Large language models \cite{touvron2023llama,jiang2023mistral,bai2023qwen}
have grown rapidly. Instruction tuning \cite{ouyang2022instructgpt} and
alignment techniques \cite{bai2022constitutional} are now standard practice.

\lipsum[3-4]

\section{Background}
\lipsum[5-6]
Retrieval-augmented generation \cite{lewis2020rag} and chain-of-thought
prompting \cite{wei2022cot} improved factual accuracy and reasoning.

\lipsum[7-8]

\section{Methods}
\lipsum[9-11]

\section{Experiments}
\lipsum[12-14]

\section{Conclusion}
\lipsum[15]

\bibliography{arxiv_refs}
\bibliographystyle{acl_natbib}

\end{document}
